\lyseisE{Επανάληψη κεφαλαίου Β4}
\solr{Β4}{1} α) ορθογώνιο παραλληλεπίπεδο \quad β) κύλινδρος \quad γ) πυραμίδα \quad δ) ορθογώνιο παραλληλεπίπεδο \quad ε) κώνος \quad στ) σφαίρα.
\solr{Β4}{2} α) Κύβος: $6$ έδρες, $12$ ακμές, $8$ κορυφές. \quad β) Τετράεδρο: $4$ έδρες, $6$ ακμές, $4$ κορυφές. \quad γ) Τριγωνικό πρίσμα: $5$ έδρες, $9$ ακμές, $6$ κορυφές.
\solr{Β4}{3} $\nu=8$: έδρες $8+2=10$, κορυφές $2\cdot8=16$, ακμές $3\cdot8=24$.
\solr{Β4}{4} $2\nu=20$, $\nu=10$: δεκαγωνική πυραμίδα με $11$ έδρες και $11$ κορυφές.
\solr{Β4}{5} Μόνο το α. Το β είναι ορθογώνιο $2\times3$ τετραγώνων: διπλώνοντας, τα τετράγωνα πέφτουν το ένα πάνω στο άλλο. Το γ δεν κλείνει, γιατί το πρώτο και το πέμπτο τετράγωνο πέφτουν στην ίδια έδρα.
\solr{Β4}{6} α) κύλινδρος \quad β) κώνος \quad γ) τετραγωνική πυραμίδα \quad δ) τριγωνικό πρίσμα.
\solr{Β4}{7} $4\cdot(5+3+2)=40$ cm.
\solr{Β4}{8} Ελεύθερη κατασκευή. Το μήκος του ορθογωνίου είναι η περιφέρεια της βάσης: $2\cdot3{,}14\cdot1{,}5=9{,}42$ cm. Το πλάτος του είναι $4$ cm και υπάρχουν δύο κύκλοι ακτίνας $1{,}5$ cm.
\solr{Β4}{9} $\mu=\dfrac{360\cdot3}{6}=180^\circ$. Ο τομέας είναι ημικύκλιο, γιατί η γενέτειρα είναι διπλάσια από την ακτίνα της βάσης.
\solr{Β4}{10} α) κύλινδρος \quad β) κώνος \quad γ) σφαίρα.
\solr{Β4}{11} Ελεύθερη κατασκευή. Μπροστά βλέπουμε ένα ορθογώνιο και από πάνω έναν κύκλο.
\solr{Β4}{12} Ακτίνα $18:2=9$ cm. Η σφαίρα δεν έχει έδρες (έχει μία καμπύλη επιφάνεια).
\solr{Β4}{13} Το εξαγωνικό πρίσμα ($6+2=8$ έδρες, $2\cdot6=12$ κορυφές). Έχει $3\cdot6=18$ ακμές.
\solr{Β4}{14} Τετραγωνική πυραμίδα: η βάση φαίνεται από πάνω ως τετράγωνο και οι ακμές προς την κορυφή ως διαγώνιοι, ενώ από μπροστά φαίνονται δύο παράπλευρες έδρες ως τρίγωνο.
\solr{Β4}{15} α) Σωστό. \quad β) Σωστό. \quad γ) Λάθος: οι παράπλευρες έδρες του πρίσματος είναι ορθογώνια. \quad δ) Σωστό.
\solr{Β4}{16} α) Λάθος: έχει μία. \quad β) Σωστό. \quad γ) Σωστό. \quad δ) Σωστό.
